\section{Storage, Load/Save and the SD
Card}\label{storage-loadsave-and-the-sd-card}

Everything you keep lives on the SD card. The card must be formatted
\textbf{FAT32}. Files are called \emph{files}, not cards.

\draftnote{This chapter summarises the filesystem and Load/Save documents, which are written for
firmware developers. Confirm every step on the device, especially the names of the Load/Save
types and the on-screen prompts.}

\subsection{Library types}\label{library-types}

\begin{paramtable}
Bank & Bank & A folder of up to 16 Scenes. Slot numbers 000 to 999; over MIDI CC0 reaches slots 0 to 127. \\
Scene & Scene & Pattern, Kit, Effect, mix, timing and chain settings for one Scene. \\
Kit & Kit & Six Instruments. \\
Instrument & Instrument & One voice patch, filed by type (Drum, Snare, Cymbal, HiHat, Alkaline). \\
Pattern & Pattern & The Pattern of a Scene (7 tracks of 128 steps), saved on its own. \\
Samples & Samples & Not a library: installs samples and wavetables into flash memory. \\
\end{paramtable}

The Effect is stored as part of its Scene. There is currently no
separate Effect library.

\verify{The Load/Save page may also offer a Settings or Globals type; the settings file is described below.}

\subsection{Load and Save}\label{load-and-save}

The Load/Save mode button opens the Load/Save pages (a second press
toggles between Load and Save).

\figplaceholder{load_save_page.png}{The Load page.}{Display close-ups: top-level Load with the type cell, slot number and name; then a nested Instrument load showing the Instrument type and slot; then the Save page with the name editor and the OK cell. Mark the busy indicator with three dots and the error overlay W followed by a code.}

\begin{itemize}
\tightlist
\item
  The \textbf{type} cell selects Bank, Scene, Kit, Instrument, Pattern
  or Samples. \textbf{Knob 1} steps through the types directly. For
  Instrument, the types offered follow the selected track's Instrument
  types.
\item
  The \textbf{slot} number and \textbf{name} show what is selected.
  Turning the encoder previews the slot; the name appears after the card
  index has been read, and stays blank while the card is busy.
\item
  Moving to a different type while an index is still being read cancels
  that read; only the settled type is read.
\item
  Writes, name updates and loading do not allow the controls to be used
  until they finish, but a watchdog ends every operation that makes no
  progress for \textbf{12 seconds}: the display then shows an error code
  beginning with \texttt{W}.
\item
  Commands are confirmed by moving to \textbf{OK} and pushing the
  encoder; \textbf{OW} appears when a Save would overwrite an existing
  slot.
\end{itemize}

\verify{The exact wording of the OK, OW and overwrite prompts.}

\subsubsection{Saving}\label{saving}

\begin{enumerate}
\def\labelenumi{\arabic{enumi}.}
\tightlist
\item
  Choose the type and slot. A slot number in use is overwritten after
  confirmation; a free slot creates a new file.
\item
  Edit the name (see below).
\item
  Move to OK and push the encoder.
\end{enumerate}

\textbf{Names.} A name stem is at most \textbf{eight characters} long. A
blank name is stored as \texttt{none.drm} for Instruments (the extension
follows the type) and as \texttt{none.fx} for the Effect. An empty file
stem becomes \texttt{inst}.

\subsubsection{Loading a Bank}\label{loading-a-bank}

A Bank load can run in the background while the sequencer plays. The
music continues from Scene 17 (a hidden copy of the Scene that was
playing); the panel is available again once the load reaches the point
where playback no longer depends on the files being read. The Load page
stays open until the load ends if the sequencer was stopped. A Bank
loaded over MIDI always replaces all 16 Scenes.

\subsubsection{Loading a Scene or Kit, Instrument and
Pattern}\label{loading-a-scene-or-kit-instrument-and-pattern}

\begin{itemize}
\tightlist
\item
  \textbf{Scene} load replaces the Scene's Pattern, Kit, Effect, mix
  settings and chain.
\item
  \textbf{Kit} load replaces all six Instruments; the Effect is not
  touched.
\item
  \textbf{Instrument} load replaces one slot. The previous state of the
  slot is kept in memory so the load can be reversed.
\item
  \textbf{Pattern} load can fan out to the Scenes in the edit mask.
\end{itemize}

\subsubsection{Morph targets}\label{morph-targets}

Morph values are stored inside Scenes, Kits and Instruments: Kits and
Instruments carry their own Morph endpoints (\texttt{{[}morph{]}}), and
a Scene stores the Morph endpoints of everything it owns. There is no
separate ``Morph sound'' file type.

\subsection{Samples}\label{samples}

\textbf{Load \textgreater{} Samples} walks \texttt{/samples} and
\texttt{/wavetables} on the card, erases the sample area of flash
memory, installs the files and writes a new catalogue.

\begin{itemize}
\tightlist
\item
  Files must be mono, 16-bit, 44.1 kHz PCM WAV.
\item
  Wavetables: one sub-folder per wavetable inside \texttt{/wavetables}.
\item
  The sample area holds five 256 KB sectors (about 1.25 MB). The
  catalogue sits at the top.
\item
  Audio is stopped during the write and restarted afterwards.
\end{itemize}

\attention{Samples are held in flash memory, not on the card. A device whose catalogue was written by
an older firmware shows STALE and must reinstall its samples before they can be used. Reinstall at
low volume before rolling the firmware back to an older version, because older firmware plays
program code as audio.}

\subsection{AutoSave}\label{autosave}

The machine saves its working state in the background: scalar settings,
each Scene's Pattern, names and the Effect. AutoSave never makes you
wait:

\begin{itemize}
\tightlist
\item
  It pauses during Load/Save, copy/clear and while a Bank load is
  running.
\item
  It restarts from the beginning after you leave the Load/Save page if
  work is still owed.
\item
  At start-up the firmware first \textbf{settles} any unfinished
  AutoSave work (this is why start-up can take a few seconds) before
  audio starts.
\item
  Settings are written to \texttt{settings.cfg} and the MIDI Learn table
  to \texttt{midilrn.cfg}.
\end{itemize}

AutoSave writes alternating hidden files so that an interrupted write
can never destroy the last good copy.

\subsection{SD card layout}\label{sd-card-layout}

\figplaceholder{sd_card_layout.png}{SD card folder layout.}{Tree view of the card root: Bank, Scene, Kit, Pattern, Instrument with a sub-folder per type (for example Alkaline containing init.alk), samples, wavetables, and the root files settings.cfg and midilrn.cfg. Mark hidden files in grey.}

\begin{paramtable}
/Bank, /Scene, /Kit, /Pattern & Library folders & Numbered folders or files with a three-digit slot and a name. \\
/Instrument & Instrument library & One sub-folder per Instrument type. \\
/samples, /wavetables & Sample sources & Used by Load > Samples. \\
settings.cfg & Settings & Global settings, read at boot. \\
midilrn.cfg & MIDI Learn & One readable line per learned CC. \\
\end{paramtable}

Hidden index and name files (\texttt{.hcindex}, \texttt{.hcnames} and
the AutoSave records) are managed by the firmware.

\textbf{Editing the card on a computer} is fine, but changes are seen
only after a power cycle; the firmware never remounts the card while
running. Do not delete the hidden files unless told to.

\subsubsection{Preparing an older card}\label{preparing-an-older-card}

A card prepared with an older firmware may have incompatible hidden
files. Delete these from the card root once and power-cycle:
\texttt{.hcnames}, \texttt{.hcnamtmp}, \texttt{.hcprms1},
\texttt{.hcprms2}.

\subsubsection{Start-up repair}\label{start-up-repair}

At start-up the firmware repairs any directories whose end-of-directory
marker was lost after deleting files. This happens once; subsequent
boots skip clean directories.

\subsection{Error codes}\label{error-codes}

\begin{paramtable}
W + code & Load/Save watchdog & An operation made no progress for 12 seconds. The code names the operation and the phase it was in. Note it and power-cycle. \\
Img BAD s:... & Flash image check & A stamped checksum in a flash sector does not match. The device waits for BAR1. See the Firmware Update chapter. \\
STALE & Sample catalogue & The sample area was written by an older firmware. Reinstall with Load > Samples. \\
\end{paramtable}

\verify{The firmware's own wording and layout for each message, and whether the flash image check stays in release firmware (the project lists this as an undecided item).}
